\section{Prospective Equal-Budget Continuation Experiment}\label{sec:prospective}
We test a separate acquisition-strategy contrast with BE/E1 and GM/E1
(19,396 parameters). Ten fresh paired seeds per teacher give forty models.
Before training, a fixed SHA256 partition holds out 20\% of sufficient
queried-slot COMMIT type buckets, common across seeds. Phase, root/context
features and distinct neighbor message types determine this key; MAX removes
multiplicity. Neither training arm supervises a held-out continuation type.
LINIAL and DETECT remain separate native targets. This is a split of input
contexts within the same intervention templates. It does not withhold all
coarser oracle-relevant configurations (quantified below).

Each pair shares 3,000 native component targets and 3,600 fresh auxiliary
rooted states. On each auxiliary state the native arm acquires one native
COMMIT target and CF one prespecified intervened COMMIT target. Both queries
must pass the training-type filter to retain that state. GM auxiliary quotas
are 1,800 ASSIGN and 1,800 REPAIR; BE uses KW only. Each model therefore
accesses 6,600 optimization targets. Native P/C conditioning is supplied to
both arms equally and charged separately: 7,200 root values plus retained
neighbor values. Simulator computation, rejected states and duplicate queries
are also recorded; this does not equate all teacher-computation costs.
The root option cycles through native-option offsets 0,1,2,3 modulo four.
Participation is the root permission, except that offset three independently
silences it with probability 1/4; neighbor messages remain native. The cursor
advances before filtering, including discarded roots. Test queries use the
same intervention and retain held-out CF types; their phase quotas are
4,000 KW or 2,000 each ASSIGN/REPAIR. This single-query design differs from
the historical four-query diagnostic. Appendix~\ref{app:prospective-method}
specifies the key, serialization, hash and acquisition order.

Adam, batch 64 and 1,500 updates are fixed before test, with 100 matched warmup
updates and the same subsequent native/auxiliary loss weighting. Paired data,
initialization, index sequences and first-stage final tensors agree. All forty
weights are frozen before generating 4,000 held-out queries per seed/teacher
on sizes 28/36 and evaluating 1,000 autonomous runs per model on sizes 32/48/64.
All exact references pass. Every final continuation has zero intersection
with its training sufficient-type set; both truth classes occur in every seed.


\begin{table}[htbp]\centering\small
\begin{tabular}{lrrlc}
\toprule Teacher & Native BA & CF BA & Paired difference [95\% interval] & Successes N / CF \\
\midrule
BE & 84.64 & 90.41 & +5.77 [-13.73,+25.26] & 9,204 / 7,405 \\
GM & 50.00 & 100.00 & +50.00 [+50.00,+50.00] & 10,000 / 10,000 \\
\bottomrule
\end{tabular}
\caption{New final experiment. BA is a percentage; differences and intervals are percentage points,
averaged over ten paired training seeds. Intervals bootstrap paired seeds
(10,000 resamples), conditional on this fixed split/protocol; individual
queries are not independent replications. Success totals are out of 10,000
autonomous runs per teacher/arm. No checkpoint selection uses these tests.}
\label{tab:prospective}
\end{table}

GM equal-type accuracy is 35.44\% native and 100.00\% CF; both arms average 26.01 native steps.

GM gains 50 percentage points in every paired seed without changing any
autonomous outcome or mean step count. Its degenerate bootstrap interval
reflects those identical observed seed differences, not certainty over all
future types or training replicas. BE's mean difference is +5.77 points, but
the interval spans zero (\(-13.73,+25.26\)). BE equal-type accuracy increases
from 51.70\% to 86.94\%, while the two endpoints weight errors differently.
CF adds false negatives in two BE seeds and reduces autonomous successes from
9,204 to 7,405; one further seed has identical failures in both arms.
Mean step counts over \emph{all} final runs (14.53 native versus 13.43 CF)
include early exits on failure. Every successful BE run takes 15 native
steps in both arms; the lower overall mean is not a successful-run speedup.
These are substantive negative results for a uniform-benefit claim.

\paragraph{Post-test diagnostic controls.}
The primary endpoint and all original outcomes above are retained.
The following controls and support audit were added after inspecting the test;
they are descriptive checks, not a revised primary analysis. The programmed
rule $G_v=E_vp_v$ uses supplied eligibility and root participation and ignores
query occupancy. It requires no training and is a diagnostic, not a learned
label-efficiency competitor.

\begin{table}[htbp]\centering\small
\begin{tabular}{llrr}
\toprule Teacher & Procedure & BA (\%) & Equal-type acc. (\%) \\
\midrule
BE & Native & 84.64 & 51.70 \\
BE & CF & 90.41 & 86.94 \\
BE & Eligibility/participation & 99.71 & 81.27 \\
GM & Native & 50.00 & 35.44 \\
GM & CF & 100.00 & 100.00 \\
GM & Eligibility/participation & 53.03 & 62.34 \\
\bottomrule
\end{tabular}
\caption{Same final queries and sufficient-type key; means over ten seeds.
The programmed diagnostic ignores occupancy. BA averages sensitivity and
specificity; equal-type accuracy first averages within each distinct
type, then across types within seed. The diagnostic is post-test.}
\label{tab:prospective-controls}
\end{table}

For BE, the programmed rule scores 99.71\% BA, above CF's 90.41\%,
whereas CF scores 86.94\% equal-type accuracy versus 81.27\%.
A single negative type accounts for 95.05--96.30\% of each 4,000-query bank.
Class balancing therefore leaves the negative subtypes highly imbalanced;
high occurrence-level BA alone does not identify the occupancy condition.
Among eligible active BE queries, native and the programmed rule accept
995/995 free options but reject 0/228 occupied options.
CF rejects 228/228 occupied options but accepts only
814/995 free options. These are pooled descriptive counts over
ten banks, not independent training replications. Appendix~\ref{app:prospective-audit}
reports the full confusion strata.

GM's native control receives 4,344--4,354 positive COMMIT labels and zero
negative COMMIT labels in ASSIGN/REPAIR per seed, while CF receives
1,038--1,098 negatives and the test has 2,420--2,491 negatives per seed.
The native model accepts all test queries, giving BA 50\%; the contrast
therefore demonstrates learning rejections absent from native color supervision.
DETECT has separate positive/negative conflict supervision and is excluded
from these color-label counts. This result does not compare CF with another
acquisition strategy supplying informative negative confirmations.

\paragraph{Support of the held-out queries.}
There are 17--23 distinct sufficient test types per BE seed and 82--92
per GM seed, despite 4,000 occurrences. For these color oracles
$O_T=E_vp_v(1-o_v)$, with $o_v$ the queried-option occupancy. A coarser
configuration $(\text{phase},E_v,p_v,o_v)$ is therefore oracle-sufficient.
CF training has already seen configurations covering 99.475--100.000\% of BE
test occurrences and 100\% of GM test occurrences. The original fine-key
holdout remains disjoint: novelty is defined by contextual combinations
rather than by this coarser configuration. There is no separate evaluation
restricted to novel logical conditions. These tests do not establish extrapolation to new logical
conditions. Per-seed support counts and logical overlap are in
Appendix~\ref{app:prospective-audit}; no split, label or model was changed.

\paragraph{Two distinct BE failure diagnoses (post-test).}
In the two deteriorated CF seeds, 105001 and 105008, COMMIT BA on the 3,600
acquired auxiliary queries is 99.510\% and 100.000\%, versus 52.756\%
and 52.094\% on the held-out contexts. Both models visited all auxiliary
queries. Failure to fit acquired labels alone does not explain these cases:
the observed COMMIT gap concerns generalization across contexts, without
identifying its optimization cause or an architectural impossibility.

Seed 105010 has a different defect: both arms complete only 204/1,000 runs,
despite final COMMIT BA 99.451\% native and 99.949\% CF. Their shared
first stage omits participation on 13/118 eligible test roots.
A replay of the original 32-case cohort localizes the first invariant
violation in one selected run: at KW step 7, vertex 31 must
recolor as the bound drops from 21 to 17, but remains at 17
because participation is false. Its proposed value 12 is free;
changing only root participation makes COMMIT accept it. Thus this state
exposes an upstream omission, rather than an incorrect veto of the actual
silent message. These are selected post-test diagnostics, without retuning,
a continued rescue rollout or an estimate of causes across all failed runs
(Appendix~\ref{app:be-failure-diagnostics}).


The primary endpoint concerns held-out continuation types; autonomous success
is secondary. This finite comparison concerns the specified acquisitions at
equal target/context access. It supports neither extrapolation to new logical
conditions nor an advantage over alternative informative-negative acquisition
strategies. BE's occurrence-level endpoint must be read with its diagnostic
baseline and type weighting; GM's control lacks negative color confirmations.
No universal CF superiority, minimum sample complexity or all-source teacher
efficiency follows.
The protocol, forty checkpoints, all seed outcomes, per-component costs and
freeze receipts accompany the new results (Appendix~\ref{app:prospective}).
